\section{The Variety of Experience}
\label{sec:experience}

If the only felt quality at the base is a tone with three values, how does experience become
so endlessly varied: the red of a sunset, the smell of coffee, the bite of fear, the click of
understanding? The answer turns on one move.

\paragraph{Charge versus quality.}
A single number can carry only a one-dimensional charge. The tone is exactly that, a point on
the single axis from pain through peace to pleasure, which is what feeling researchers call
\emph{valence}, how good or bad something feels. But quality, the specific character of an
experience, is not one-dimensional. Red and middle C and the smell of coffee are wildly
different in a high-dimensional way, so quality cannot live in the tone. It must live in the
\emph{structure}: the pattern of tones across the mesh, where they sit, how they move, how
they connect.

\begin{principle}[Charge and quality]
The tone is the charge (valence), the only thing that is intrinsically felt as good, neutral,
or bad. Everything else, the specific quality, is structure. Every experience is a structured
distribution of ternary charge across the mesh.
\end{principle}

This is why a tiny alphabet suffices. Twenty-six letters give all of literature, four DNA
bases all of life, two bits all of computation. The variety is never in the alphabet, always
in the arrangement. It also explains a felt fact: an emotion feels hot and saturated while a
thought feels cool and clear, because an emotion is dense in tone and a thought is sparse in
tone but rich in structure.

\paragraph{The dimensions of variety.}
Seven degrees of freedom generate the range. Which subsystem the pattern lives in sets the
\emph{modality} (sight, sound, thought). The configuration of tones within it is the specific
\emph{content}. The density and rate is the \emph{intensity} (vividness, loudness, arousal).
The tone balance is the \emph{valence}. The trajectory over beats is the \emph{dynamics}
(motion, melody, the arc of a feeling). The connections to other patterns are the
\emph{binding} (meaning, association, object-hood). And the height in the recursion is the
\emph{level}, low for raw sensation, higher for emotion and thought.

\paragraph{The senses.}
The proposal is that a sense can be thought of as a subsystem with its own geometry, and that
the geometry is the shape of that sense's quality-space. Sight might be considered a roughly
two-dimensional sheet, so that a visual experience would be a pattern over a sheet, with color
a small extra dimension at each point. Hearing might be taken as roughly one-dimensional, laid
out by frequency and carrying strong dynamics, so that pitch is a position and a melody a
trajectory. Touch could be a sheet laid out like the body's surface. Taste and smell would be
high-dimensional chemical spaces, and smell's space may be so large and unordered that smells
are vivid yet hard to name, which is the kind of thing this picture would anticipate. The
suggestion, then, is that two senses feel different because their sub-meshes are differently
shaped in different places, the difference in quality tracking a difference in structure. These
are working analogies meant to make the idea imaginable, not settled assignments.

\paragraph{Emotions and thought.}
In the same spirit, emotions can perhaps be pictured as higher-level patterns, characteristic
distributions of pain, peace, and pleasure across the self, each with a level of arousal, a pull
to act (the urge), and a target. Joy would be pleasure-rich, open, and broad. Fear would
be pain-rich, high in arousal, withdrawing from a threat. Anger would be an approach toward an
obstacle, sadness a withdrawal from a loss, and the social emotions such as guilt, pride, and
awe would be aimed at higher-level targets like the self-model or a much larger whole. Thought might be activity in
the abstract upper levels, where a concept is a stable pattern (an attractor), thinking is a
trajectory that moves among them, reasoning is a kind of computation, and inner speech is the
reconstruction of active concepts into words. The general claim is the firm one, that three
tones can become the whole texture of mind because the variety lives in the structure rather
than the alphabet. The specific assignments are offered as plausible illustrations that may well
be refined.
